\documentclass[a4paper,12pt]{article}

%-------------------------------------%
%-------------------------------------%
%                                     %
%        parcial din rotació 20 21    %
%         Física  1r btx  bio         %
%                                     %
%-------------------------------------%
\usepackage[utf8]{inputenc}
\usepackage[catalan]{babel}
\usepackage{amssymb}
\usepackage[symbol]{footmisc}
\usepackage{multicol}
\usepackage{amsmath}
\usepackage{cancel}
\usepackage{xcolor}
\usepackage{graphicx}
\usepackage[most]{tcolorbox}
\usepackage{tikz}
\usepackage{pgfplots}
\pgfplotsset{compat=1.18}
\usepackage{eurosym}
\usepackage{wasysym}
\usetikzlibrary{decorations.pathmorphing}
\usetikzlibrary{arrows,positioning,shapes,fit,calc}
\usepackage{hyperref}
\hypersetup{hidelinks,}
\usepackage{fancyhdr}
\pagestyle{fancy}
\fancyhf{}
\rhead{\scriptsize A.Arroyo (aarroyo+batx@stjosep.org)}
\lhead{\scriptsize Parcial Dinàmica rotació. Física 1r Batx. 2025-26}
\lfoot{\href{http://creativecommons.org/licenses/by-nc-nd/4.0/}{\includegraphics[width=1.5cm]{./../../../cc}}}
\rfoot{\thepage}
\setlength\headheight{14pt}
\begin{document}
%\begin{enumerate}
%\item
%\begin{enumerate}
%\item
\noindent\textbf{1.} (a) 
Tal com planteja l'enunciat, es tracta d'un exemple de corba peraltada amb velocitat màxima, llavors
\begin{tikzpicture}
\draw node[left,text width=4cm, align=justify, rounded corners,fill=red!10,inner sep=2ex]
{\emph{\underline{Velocitat màxima}}};
\end{tikzpicture}


\begin{tikzpicture}
\draw[step=1cm,white,thin] (-6,-4) grid (6,3);
%\draw[step=0.1cm,blue!10,ultra thin] (-6,-4) grid (6,3);
\filldraw [brown!30!black](0,0) circle (0.1cm);
\draw[thick](-3,2)--(2,-3)--(-3,-3)--cycle;
\shadedraw[xshift=0.2cm,yshift=0.2cm,rotate=-45,left color=white, right color= {rgb:red,2;green,2;blue,3}, draw= blue!75!black, thick] (-0.7,-0.7) rectangle (0.7,0.5);
\filldraw[gray!40,draw= blue!75!black](-0.85,-0.15)--(-0.65,0.05)--(-0.5,-0.1)--(-0.7,-0.3)--cycle;
\filldraw[gray!40,draw= blue!75!black](-0.15,-0.85)--(0.05,-0.65)--(-0.1,-0.5)--(-0.3,-0.7)--cycle;
\draw [black!60!green,thick,->,>=stealth](0,0) -- (0,-2.1);
\draw [black!70!yellow,thick,->,>=stealth](0,0) -- (2,2);
\draw [thick,->,>=stealth](0,0) -- (1.5,-1.5);
\draw (1.2,-1.7) node[scale=0.8] {$F_{f}$};
\draw (1.5,1) node[scale=0.8] {$N$};
\draw (-0.5,-2) node[scale=0.8] {$mg$};
\draw [densely dashed](-2.5,0)--(2.5,0);
\draw [densely dashed](0,-2.5)--(0,2.5);
\draw [thick,->,>=stealth](1.8,0) -- (2,0);
\draw [thick,->,>=stealth](0,1.8) -- (0,2);
\draw [thick,->,>=stealth](0,-1.3) -- (0,-1.5);
\draw [thick,->,>=stealth](1.3,0) -- (1.5,0);
\draw [loosely dashed](1.5,0)--(1.5,-1.5)--(0,-1.5);
\draw [loosely dashed](0,2)--(2,2)--(2,0);
\filldraw[fill=green!20!white, draw=green!50!black]
(2,-3) -- (1.5,-2.5) arc (135:187:6mm) -- cycle;
\draw(0.16,0.18) arc (45:90:0.2cm);
\draw (1,-2.75) node[scale=0.8] {$\alpha$};
%\draw(0,0) arc (-45:-90:0.3cm);
\draw(0.5,0) arc (0:-45:0.5cm);
\draw (1.9,-0.33) node[scale=0.6] {$N_{x}$};
\draw (0.2,1.7) node[scale=0.6] {$N_{y}$};
\draw (-0.3,-1.3) node[scale=0.5] {$F_{f_{y}}$};
\draw (1.2,-0.3) node[scale=0.5] {$F_{f_{x}}$};
\draw (0.15,0.3) node[scale=0.5] {$\alpha$};
\draw (0.6,-0.2) node[scale=0.4] {$\alpha$};
\end{tikzpicture}






-----------------------------------
\begin{center}
\begin{tikzpicture}
\draw(-3,0)arc(180:265:3);
\draw(3,0)arc(0:-85:3);
\draw[very thick,cyan!70!green](-0.26,-3)--(0.26,-3);
\draw(-0.25,-2.9)rectangle(0.25,-2.5);
\draw(-3,0)rectangle(-2.6,-0.5);
\draw[->,dashed](0,0)--(-2.12,-2.12);
\draw[->](0,-2.75)--+(0,1);
\draw[->](0,-2.75)--+(0,-1);
\filldraw(0,-2.7)circle(0.025);
\filldraw(0,0)circle(0.025);
\node[scale=0.75] at (0.5,-3.75) {$mg$};
\node[scale=0.75] at (0.5,-1.75) {$N$};
\node[scale=0.75] at (-0.75,-0.25) {$R$};
\end{tikzpicture}
\end{center}
%\item

\noindent (b) En el moment en que passa per sobre l'esquerda podem escriure
$$N-mg=m\frac{v^{2}}{R}$$
si no volem superar el valor màxim de la normal, la velocitat ha de ser
$$v_{max}=\sqrt{R\cdot\frac{N_{max}-mg}{m}}=\sqrt{15\cdot\frac{1000-80\cdot9,8}{80}}=6,364\,m/s$$
%\end{enumerate}
%\item
\noindent\textbf{2.} Fet als \textcolor{cyan!20!green}{\href{https://artur-sjo.github.io/Fisica_1r_Btx/Fis_1r.pdf}{apunts del curs}}, secció 4.10 \emph{La corba peraltada}, pàgines 82, 83.
\bigskip

%\item
%\begin{enumerate}
%\item
\noindent\textbf{3.} (a) Fem un diagrama vist des de dalt i de costat
\begin{center}
\begin{tikzpicture}
\draw(0,0)circle(2);
\filldraw[white,draw=black](1.5,-0.5)rectangle(2.5,0.5);
\draw(0,0)--(2,0);
\filldraw(0,0)circle(0.05);
\draw[->](2,0)--(1,0);
\node[scale=0.75] at (-0.75,-0.25) {$R$};
\draw[->,dashed](0,0)--(-1.414,-1.414);
\draw(5,0)--(7,0);
\filldraw(5,0)circle(0.05);
\draw(6.5,-0.5)rectangle(7.5,0.5);
\draw[->](7,0)--+(-1,0);
\draw[->](7,0)--+(0,-1.5);
\node[scale=0.75] at (1,0.25) {$T$};
\node[scale=0.75] at (6,0.25) {$T$};
\node[scale=0.75] at (7,-1.75) {$mg$};
\end{tikzpicture}
\end{center}
Com es pot veure, el pes i la tensió estan desacoblats, i nomes aquesta darrera contribueix a la força centrípeta.
%\item
\noindent(b) Ara, podem escriure
$$T=m\omega^{2}R$$
d'on
$$\omega=\sqrt{\frac{T}{Rm}}=\sqrt{\frac{300}{1\cdot3}}=10\,rad/s$$
%\end{enumerate}
%\item
\noindent\textbf{4.} Al girar, la superfície del tambor de la rentadora proporciona la força centrípeta a la roba. Al haver-hi uns petits orificis, l'aigua es pot escolar per ells (però no la roba). És a dir, l'aigua no surt del tambor per gravetat, si no perquè no té suport per girar dins el tambor. D'aquesta manera la roba s'asseca una mica. 
%\end{enumerate}
\end{document}
